\section{Task 1: Multimodal Classification of Product Defects}
\label{sec:t1}
% Code: kode/task1_klasifikasi/ (python run_all.py --task 1).
% Fill [TBD] entries from kode/results/task1/experiment_summary.json,
% per_class_f1_comparison.csv, and likely_mislabels.csv, then uncomment the figure blocks.

We classify 16{,}000 test product units into seven defect classes (\texttt{normal}, \texttt{gores},
\texttt{noda}, \texttt{retak}, \texttt{gelembung}, \texttt{las\_dingin}, and \texttt{kontaminasi})
from 48{,}000 labelled training units. Each unit combines a $64 \times 64$ grayscale surface image
with a $3 \times 128$ press sensor recording (temperature, vibration, current). Evaluation on the
hidden test set uses macro-F1 over all seven classes.

\subsection{Data exploration}
\label{sec:t1-eda}

\paragraph{Class imbalance and metadata.} Nearly half of the 48{,}000 training units are
\texttt{normal} (23{,}888 units, 49.8\%), while the six defect classes range from 11.7\%
(\texttt{gores}, 5{,}623) down to 4.9\% (\texttt{kontaminasi}, 2{,}357); Table~\ref{tab:t1-dist}
lists the full distribution. Units come from 12 press machines, 150 raw-material batches, and three
balanced shifts (\texttt{pagi}, \texttt{siang}, \texttt{malam}).

\paragraph{Missing sensors and train--test domain shift.} In the training set, 2{,}352 units
(4.90\%) have completely missing sensor recordings (all \texttt{NaN}) and 3{,}863 units (8.05\%)
contain at least one \texttt{NaN}; the test set shows almost identical rates (813 units or 5.08\%
completely missing; 1{,}298 or 8.11\% with any \texttt{NaN}). In contrast, the production-line
composition shifts sharply between train and test: training data is dominated by Line~A
(35{,}964 units, 74.9\% vs.\ 12{,}036 or 25.1\% on Line~B), whereas the test set is dominated by
Line~B (9{,}607 units, 60.0\% vs.\ 6{,}393 or 40.0\% on Line~A).

\begin{table}[t]
  \centering
  \caption{Class distribution in the 48{,}000 training units and the 16{,}000 submitted test
  predictions (\texttt{prediksi\_klasifikasi.csv}).}
  \label{tab:t1-dist}
  \small
  \begin{tabular}{lrrrr}
    \toprule
    Class & Train count & Train share & Test pred.\ count & Test pred.\ share \\
    \midrule
    \texttt{normal}      & 23{,}888 & 49.8\% & 9{,}619 & 60.1\% \\
    \texttt{gores}       &  5{,}623 & 11.7\% & 1{,}240 &  7.8\% \\
    \texttt{noda}        &  4{,}773 &  9.9\% & 1{,}332 &  8.3\% \\
    \texttt{las\_dingin} &  4{,}323 &  9.0\% & 1{,}326 &  8.3\% \\
    \texttt{retak}       &  3{,}606 &  7.5\% &   806   &  5.0\% \\
    \texttt{gelembung}   &  3{,}430 &  7.1\% &   960   &  6.0\% \\
    \texttt{kontaminasi} &  2{,}357 &  4.9\% &   717   &  4.5\% \\
    \midrule
    Total                & 48{,}000 & 100.0\% & 16{,}000 & 100.0\% \\
    \bottomrule
  \end{tabular}
\end{table}

\subsection{Preprocessing and architectures}
\label{sec:t1-arch}

\paragraph{Preprocessing.} Training data is split 80/20 with class stratification
(\texttt{random\_state=42}), and all normalization statistics are fitted on the 38{,}400 training
samples only. Because raw sensor arrays are stored in \texttt{float16}, valid (non-\texttt{NaN})
readings are cast to \texttt{float64} before computing per-channel means and standard deviations;
computing statistics directly in \texttt{float16} overflows on 4.9M entries and zeroes out the sensor
branch. Missing entries are filled with 0.0 after standardization. During training we apply random
horizontal/vertical flips to images and modality dropout ($p=0.15$, zeroing the entire sensor
tensor) so the model learns to fall back on the image when sensors are absent.

\paragraph{Encoders and fusion strategies.} Both modalities use compact residual networks~\cite{he2016}
producing 128-dimensional representations:
\begin{itemize}
  \item \textbf{Image encoder (2D ResNet):} initial $3\times 3$ convolution (16 channels) followed by
  four strided residual stages (32, 64, 128, 128 channels) reducing $64\times 64$ to $4\times 4$, then
  global average pooling and a linear projection to $\mathbb{R}^{128}$ (or 16 spatial tokens of dim
  128 for attention).
  \item \textbf{Signal encoder (1D ResNet):} initial kernel-5 convolution (16 channels) followed by
  four strided 1D residual stages (32, 64, 128, 128 channels) reducing length 128 to 8, then global
  average pooling and projection to $\mathbb{R}^{128}$ (or 8 temporal tokens for attention).
  \item \textbf{Late fusion:} concatenates the two 128-d embeddings and passes them through a
  two-layer MLP ($256 \to 128 \to 7$) with BatchNorm, ReLU, and dropout $0.3$.
  \item \textbf{Cross-attention fusion:} multi-head attention~\cite{vaswani2017} (4 heads,
  $d=128$) where the 16 image spatial tokens query the 8 signal temporal tokens, followed by residual
  LayerNorm, mean pooling of both streams, concatenation, and an MLP classifier.
\end{itemize}

\subsection{Modality value, label noise, domain shift, and missing sensors}
\label{sec:t1-experiments}

All comparison experiments train for 15 epochs with AdamW ($\text{lr}=10^{-3}$, weight decay
$10^{-4}$), cosine annealing, batch size 128, and seed 42; the final ensemble trains for 18 epochs
across three seeds (42, 123, 999). Before full training, a sanity check overfits one 32-sample batch
for 40 epochs (\texttt{train\_eval.py --sanity-check}), verifying cross-entropy loss $< 0.1$ and
macro-F1 $= 1.00$.

\begin{table}[t]
  \centering
  \caption{Task 1 validation results on the stratified 20\% split (9{,}600 units). Values marked
  \textbf{[TBD]} are populated directly from \texttt{kode/results/task1/experiment\_summary.json}.}
  \label{tab:t1-results}
  \small
  \begin{tabular}{lr}
    \toprule
    Experiment / Model & Val Macro-F1 \\
    \midrule
    Image-Only 2D ResNet                                & \textbf{[TBD]} \\
    Sensor-Only 1D ResNet                               & \textbf{[TBD]} \\
    Multimodal Late Fusion (global norm, LS $\epsilon=0.1$) & \textbf{[TBD]} \\
    Multimodal Cross-Attention (global norm, LS $\epsilon=0.1$) & \textbf{[TBD]} \\
    Late Fusion without Label Smoothing ($\epsilon=0.0$) & \textbf{[TBD]} \\
    Late Fusion + Per-Line Norm (image \& sensor)       & \textbf{[TBD]} \\
    Cross-Attention + Per-Line Norm (image \& sensor)   & \textbf{[TBD]} \\
    Best Model with 100\% Missing Sensor                & \textbf{[TBD]} \\
    Final 3-Seed Mean $\pm$ Std (seeds 42, 123, 999)    & \textbf{[TBD]} \\
    Final 3-Seed Probability Ensemble                   & \textbf{[TBD]} \\
    \bottomrule
  \end{tabular}
\end{table}

\paragraph{Value of each modality.} Comparing \texttt{ImageOnlyClassifier} and
\texttt{SignalOnlyClassifier} per class (\texttt{per\_class\_f1\_comparison.csv}) isolates which
defects are visual (surface defects such as \texttt{gores}, \texttt{noda}, \texttt{retak}) versus
sensor-driven (such as \texttt{las\_dingin}), and shows why multimodal fusion outperforms either
modality alone. \textbf{[TBD: insert key per-class F1 numbers and uncomment
\texttt{t1\_modality\_value\_comparison.png} once added to \texttt{laporan/figures/}.]}

% \begin{figure}[t]
%   \centering
%   \includegraphics[width=\linewidth]{t1_modality_value_comparison}
%   \caption{Per-class validation F1 comparing image-only, sensor-only, late fusion, and cross-attention.}
%   \label{fig:t1-modality}
% \end{figure}

\paragraph{Imperfect labels.} Because labels come from factory inspectors, we flag likely mislabels
on validation where the trained multimodal model disagrees with the inspector label at confidence
$\ge 0.85$ (\texttt{detect\_likely\_mislabels} in \texttt{run\_experiments.py}, saved to
\texttt{likely\_mislabels.csv}), and compare standard cross-entropy ($\epsilon=0.0$) against label
smoothing~\cite{szegedy2016} ($\epsilon=0.1$). \textbf{[TBD: insert flagged mislabel count/share and
label-smoothing delta from \texttt{experiment\_summary.json}.]}

\paragraph{Domain shift across production lines.} To address the Line~A $\to$ Line~B shift (75/25\%
in train vs.\ 40/60\% in test), we evaluate macro-F1 separately on Line~A and Line~B under global
normalization versus per-line standardization of both image pixels and sensor channels (using
training-split per-line statistics). \textbf{[TBD: insert Line~A and Line~B macro-F1 before and after
per-line normalization.]}

\paragraph{Missing modality robustness.} Trained with 15\% sensor modality dropout, the chosen
multimodal model is evaluated both on naturally missing-sensor validation units and under a 100\%
missing-sensor stress test (all sensors set to \texttt{NaN}), confirming sensible image-only
fallback. \textbf{[TBD: insert natural and 100\%-missing macro-F1.]}

\subsection{Final model and test predictions}
\label{sec:t1-final}

The best-performing architecture and normalization setting on validation is trained across three
random seeds (42, 123, 999) for 18 epochs, and test predictions average the softmax probabilities of
all three checkpoints. The resulting \texttt{prediksi\_klasifikasi.csv} contains all 16{,}000 test
units in exact ID order with all seven classes represented in proportions consistent with the
training distribution and the Line~A $\to$ Line~B shift (Table~\ref{tab:t1-dist}).
\textbf{[TBD: uncomment learning-curve and confusion-matrix figures once added to
\texttt{laporan/figures/}.]}

% \begin{figure}[t]
%   \centering
%   \includegraphics[width=0.95\linewidth]{t1_confusion_matrix_final}
%   \caption{Normalized validation confusion matrix of the final 3-seed ensemble.}
%   \label{fig:t1-cm}
% \end{figure}
